\chapter{Conclusion}
\label{ch:conclusion}

\section{Summary of Contributions}
\label{sec:summary}

This thesis studied key--value pair extraction from business documents with a
generative baseline and a layout-aware discriminative system. Its main
contributions are:

\begin{enumerate}
  \item \textbf{Dataset characterisation:} The study characterises 9{,}124 unique
  coordinate-bearing KVP10k training pages with two weakly separated
  annotation-geometry clusters and a separate spatial-link audit.

  \item \textbf{Discriminative model:} The study adapts a pretrained
  LayoutLMv3-base encoder with a token classifier and a span-level relation
  linker. The final V4 model uses text and bounding boxes with constant blank
  visual input, so it does not use document-specific image information.

  \item \textbf{Controlled KVP10k evaluation:} The study compares the final corrected V4 with a reconstructed Mistral-7B
  baseline on the same 581 usable test pages. V4 reaches Regular text+location
  F1 of $0.345$, compared with $0.662$ for the reconstructed Mistral baseline on the same prepared 581-page test subset.

  \item \textbf{Diagnostic analysis:} The token-to-span progression, corrected
  entity and pair metrics, cluster evaluation, and qualitative examples suggest
  that relation linking remains a major limitation of the current pipeline.
  Because several corrections were applied together, the study does not assign
  an independent causal effect to linking granularity.
\end{enumerate}

The dataset and geometry analysis, implementation and evaluation of the
discriminative LayoutLMv3 pipeline, diagnostic studies, and reported V4
analyses were carried out as part of this thesis; KVP10k, its released
benchmark, pretrained LayoutLMv3, and the published Mistral approach are prior
resources and methods used as foundations and comparison points.

Stage~4a remains a methodological entity-pre-training phase. No reliable
checkpoint survives, and its historical metric did not penalise KEY--VALUE
confusion correctly. It therefore has no reportable numerical result.

\section{Research Questions Revisited}
\label{sec:rq_revisited}

\subsection{RQ1: Layout Characteristics}

Page-level clustering identifies two weakly separated annotation-geometry
regimes. Cluster~0 has a higher box count and broader spatial spread, while
Cluster~1 has fewer, larger boxes in a more compact region. Two principal
components explain 52.8\% of the standardised feature variance, and the
silhouette score is 0.222. These groups should not be interpreted as document genres or text-density
classes because the clustering is based on annotation geometry. In a separate exploratory audit, linked pairs have a mean normalised
centre distance of 0.120. This result suggests that spatial relations can be
useful, but it does not establish a causal requirement.

\subsection{RQ2: Discriminative vs.\ Generative}

The final V4 system reaches Regular text+location F1 of $0.345$, while the
reconstructed Mistral baseline reaches $0.662$ on the same prepared 581-page
test subset. The reconstructed value is also numerically above the published
Mistral result of $0.611$, but differences in test coverage, native-text
extraction, preprocessing, and implementation prevent a direct superiority
claim. V4 therefore remains below the generative baseline in the controlled
internal comparison.

\subsection{RQ3: Linking Granularity}

V1 uses dense token-level relation candidates, which creates a substantially
larger candidate space than the later span-level formulation. For illustration,
100 candidate key tokens and 100 candidate value tokens with about 20 positive
relations would produce about 10{,}000 candidate pairs and an imbalance of
roughly $500{:}1$. In the illustrative span-level case, 15 positive links among
225 candidates give 210 negatives, or approximately $14{:}1$. The V2 results are difficult to interpret because the implementation contained
a mismatch between predicted spans and ground-truth relation supervision, an
over-strict bounding-box filter, and cross-product link labels. Final V4 uses
ground-truth entity spans during training and predicted spans during inference,
and corrects the bounding-box and link-label pipeline issues. It reaches
official Regular text+location F1 of $0.345$. Because several procedures change together, the
experiment does not identify the independent effect of one correction or of
span-level linking alone.

\subsection{RQ4: Success and Failure Factors}

V4 reaches Regular text+location F1 of $0.305$ in Cluster~0 and $0.410$ in
Cluster~1. Mistral has the opposite order, with $0.679$ and $0.633$. This is an
association with annotation geometry, not evidence about semantic document
types. Qualitative cases show correct local links, plausible links to
neighbouring values, missed short or abbreviated values, and KEY--VALUE
confusion. These examples are illustrative and are not statistical proof.

The corrected V4 token-level entity scores are numerically higher than the
complete-pair score, although these metrics are not directly comparable.
Together with the observed wrong-link examples, this suggests that relation
linking remains a major internal limitation of the current pipeline. Location-only evaluation is less
restrictive than joint text+location evaluation, so its higher score does not
mean that text reduces model performance.

\section{Practical Implications}
\label{sec:implications}

The direct V4 model emits Regular pairs. The unchanged recovery method adds
unkeyed and unvalued outputs and raises All text+location F1 from $0.229$ to
$0.261$. It does not change Regular F1. This provides a low-cost coverage
extension, but its special-category accuracy remains limited.

The results suggest that relation modeling is an important priority for further
development. Useful future directions include constrained
one-to-one matching, hard-negative relation training, explicit row and column
features, learned unkeyed and unvalued classes, confidence calibration, better
handling of short values and long documents, and a controlled experiment with
real document images.

\section{Closing Remarks}
\label{sec:closing}

V4 is the final corrected LayoutLMv3 experiment in this thesis. It provides
a controlled result for text-and-layout processing without genuine image
content. Its official Regular text+location F1 of $0.345$ remains below the
reconstructed Mistral result. The main diagnostic finding is that correct
entity predictions do not guarantee correct key--value links. The observed
wrong-link cases motivate further work on relation modelling and controlled
evaluation with genuine visual input.
